This chapter surveys the prior work the four contributions build on,
in seven strands: the transformer architecture and the per-token
cost compression targets (Section~\ref{sec:rw_transformer}); the
scaling-law trajectory (Section~\ref{sec:rw_scaling}); the published
energy and economic measurements of inference at scale
(Section~\ref{sec:rw_cost}); the training-free compression literature
(Section~\ref{sec:rw_compression}); multi-agent systems and agentic
memory (Section~\ref{sec:rw_multiagent}); retrieval-augmented
generation and long-context benchmarks (Section~\ref{sec:rw_rag});
and privacy-aware retrieval (Section~\ref{sec:rw_privacy}). The
consolidated gap statement is Section~\ref{sec:rw_gap}.

\section{The Transformer Architecture and the Cost of Context}
\label{sec:rw_transformer}

The large language models this thesis builds on are all
transformers~\cite{vaswani2017attention}. The transformer's defining
mechanism, \emph{self-attention}, computes $N^2$ pairwise
compatibility scores
for a sequence of length $N$, so prefill compute scales as
$L\cdot(N^2 d + N d^2)$ over $L$ blocks of hidden size $d$, and the
key-value cache that decoding reads back grows linearly in $N$. The
planners this thesis evaluates are decoder-only models (Qwen, Llama,
DeepSeek, Phi-3), whose prefill cost $\approx 2PN$ (parameter count
$P$, context length $N$) halves when the context halves: the lever
context compression operates on.

Scaling laws~\cite{kaplan2020scalinglaws, hoffmann2022chinchilla}
gave the field a recipe to spend \emph{training} compute productively
but no corresponding recipe for \emph{inference} compute, which is
paid per query for the model's deployment lifetime and dominates
training cost at production scale. Context compression is to
inference cost what scaling laws are to training cost: a structural
lever that converts a linear cost into a constant one.

\section{The Scaling Era of LLMs}
\label{sec:rw_scaling}

The trajectory from GPT-2 (1.5B, 2019) to GPT-3 (175B) to the
2025--2026 frontier (Llama-3, Qwen-2.5, DeepSeek V3/V4, Claude
Sonnet 4.x, GPT-oss) marks four orders of magnitude of growth in
deployed model size, each step raising per-query cost and the demand
for context-engineering that controls how much input the model
reads. The two-to-three-order gap between frontier-cloud and
on-premises per-token cost is what makes compression an operational
concern. Context windows grew in parallel (2K to multi-million
tokens), but a model's effective use of them lags the nominal
size (the ``lost in the middle'' effect~\cite{liu2024lostinmiddle,
hsieh2024ruler}), so tokens the model will not use well are pure
cost. The model-scale trend is exactly what Corollary~1 of
\charef{experiments} tests: whether the cliff position depends on
where on the scale curve the planner sits.

\section{The Cost of Running LLMs at Scale}
\label{sec:rw_cost}

The motivation chapter frames the work around per-token cost; this
section anchors that framing in published measurements.

\paragraph{Energy.}
Luccioni \emph{et~al.}'s \emph{Power Hungry
Processing}~\cite{luccioni2024powerhungry} measures inference
wall-power across open-weight models (on the order of $0.1$~Wh per
inference for a 6--7B model, scaling near-linearly with size) and
separates a per-request fixed cost from a per-token cost that
dominates for long prompts. Samsi \emph{et~al.}~\cite{samsi2023words}
and Pope \emph{et~al.}~\cite{pope2023efficiently} corroborate that
energy and serving cost grow with the number of tokens processed
once the prefill dominates: the long-prompt regime compression cuts.

\paragraph{Token economics.}
At $\approx$~USD~3 / 15 per million input / output tokens on a
representative frontier tier (Anthropic Claude Sonnet; cheaper tiers
such as GPT-4o-mini run $\approx20\times$ lower), one million
$10{,}000$-token calls per day is
USD~$30{,}000$/day of input cost; halving the prompt saves half of
that. On-premises marginal cost is two orders lower per token but the
same lever applies, and the politeness tokens (``Hello'', ``Thanks'',
restated history) that opened \charef{introduction} convert directly
to kilowatt-hours regardless of whether they change the
answer~\cite{altman2025pleasethankyou}. The TalentAdore production
accounts confirm this pricing.

\paragraph{Multi-agent multiplier.}
Anthropic's multi-agent report~\cite{anthropic2025multiagent}
observes that token usage explains $\approx 80\%$ of the performance
variance they see, much of it shared context flowing between agents;
the follow-on context-editing report~\cite{anthropic2025contextediting}
reports an $84\%$ token reduction on a 100-turn web-search
evaluation. These establish that the cost-of-context question is
operationally pressing for production multi-agent deployments; what
is missing, and what \charef{experiments} supplies, is the
structural answer to \emph{how aggressively compression can be
applied before coordination breaks}.

\section{Context Compression: the Breakthrough Works}
\label{sec:rw_compression}

The training-free context-compression literature has matured in
the past three years into four recognisable families, surveyed
comprehensively by Li \emph{et~al.}~\cite{li2025promptsurvey}; the
most aggressive recent variants push compression to the extreme of
a single-token document representation
(xRAG~\cite{cheng2024xrag}). This
subsection surveys the four families, with a comparison table at the end
that maps the families against the design dimensions that matter
for multi-fragment evaluation. Each family closes with a
``what is missing in prior work'' line that the rest of the
thesis addresses.

\subsection{Token-level pruning (the LLMLingua family)}

The first context-compression technique to be widely adopted was
LLMLingua~\cite{jiang2023llmlingua}: a small autoregressive language
model scores tokens by their per-step perplexity contribution and
drops those whose removal least disturbs the conditional distribution. LongLLMLingua
\cite{jiang2024longllmlingua} extends the coarse-to-fine ranking
with a question-aware step to handle long retrieval-augmented
contexts. LLMLingua-2~\cite{pan2024llmlingua2} trains an
XLM-RoBERTa token classifier on data distilled from a more
capable model, producing a task-agnostic compressor that is
faster at inference than the original LLMLingua because it does
not require running a generative model at compression time. The
LLMLingua-2 classifier is the principal hard-prompt compressor of
this thesis; its evaluation centres on NarrativeQA F1
within $2$~pp and HotpotQA F1 within $3$~pp of reported
numbers~\cite{kocisky2018narrativeqa, yang2018hotpotqa} as the
calibration target.

\emph{What is missing.} All three LLMLingua papers report
single-agent QA-F1 retention; none measures whether the surviving
content lets a planner combine information across multiple compressed
fragments. H1 of \charef{experiments} measures exactly that.

\subsection{Selective and instruction-aware compression}

Li \emph{et~al.}'s \emph{Selective
Context}~\cite{li2023selective} prunes tokens
with low self-information under the model's own conditional
distribution; the technique requires no training but does
require running a scoring forward-pass at compression time. Xu
\emph{et~al.}'s \emph{RECOMP}~\cite{xu2024recomp} compresses long retrieval contexts
with a small task-aware extractor trained to summarise the
retrieved passages for the downstream question. The
instruction-aware filter of this thesis sits in the same family
but is deliberately simpler: a TF-IDF pruning step, then a
cross-encoder re-ranker (\texttt{BAAI/bge-reranker-base}), then a
stop-word and short-token trim, with no LM in the loop. This gives a
deterministic, training-free, low-cost baseline that the more
sophisticated compressors must beat.

\emph{What is missing.} This literature's task-awareness is
question-conditioned, but in a multi-fragment workflow the
fragment-level task is to preserve information a yet-unknown planner
will combine with other fragments, the gap the instruction-aware
filter and the multi-fragment evaluation address.

\subsection{Learned compression (parameter-cost compressors)}

Mu \emph{et~al.}~\cite{mu2023gist} introduced \emph{gist tokens},
special positions trained with a masked attention pattern so that
the model learns to compress instruction prefixes into a small
number of virtual tokens. Chevalier
\emph{et~al.}~\cite{chevalier2023autocompressor} extended this
with \emph{AutoCompressor}, which fine-tunes a language model to
process long documents segment by segment and produce summary
vectors. Ge \emph{et~al.}~\cite{ge2024icae} formalised the
architecture as the \emph{In-context Autoencoder (ICAE)}, which
trains a LoRA-adapted encoder to produce a fixed number of
memory tokens from an input context. Wang
\emph{et~al.}~\cite{wang2024icf} present the In-Context Former,
a faster cross-attention variant. Rae
\emph{et~al.}~\cite{rae2020compressive} introduced the
Compressive Transformer architecture, which is the closest
single-context analogue of the cliff bound this thesis derives.

\emph{What is missing.} Learned compression requires training,
introducing a training-budget asymmetry that muddies
cross-compressor comparison. This thesis stays training-free so the
cliff comparison turns on the algorithm rather than the training
budget; re-evaluating learned compressors under the same methodology
is named as follow-on work.

\subsection{Extractive small-LLM span selection}

Recent work uses small instruction-tuned LLMs as extractive
compressors that produce verbatim spans of the source rather than
abstractive summaries. The candidate small LLMs are the Phi-3
family~\cite{phi3} and other $\sim 3$B-parameter instruction-tuned
models. Extractive compression is privacy-friendly because the
compressed output is by construction a subset of the source
tokens; this is one of the reasons it appears in this thesis as
a compressor option.

\emph{What is missing.} Published guidance on enforcing the
extractive constraint reliably is thin; small LLMs insert
function-word filler even when told not to. This thesis's Phi-3-Mini
extractive compressor ships a verifier with a $15\%$ novel-token
tolerance and an LLMLingua-2 fallback (Section~\ref{sec:compressors}).

\subsection{Comparison table}

\begin{table}[htbp]
\centering
\caption{The compressor families this thesis surveys, along the
design dimensions that matter for multi-fragment evaluation. The
third column is the property H1 of \charef{experiments}
turns on: no prior compressor was evaluated on a controlled
multi-fragment coordination sweep before this thesis.}
\label{tab:rw_comparison}
\small
\begin{tabular}{p{4.5cm}p{2cm}p{2cm}p{2cm}p{2.5cm}}
\hline
Family & Training\hspace*{0pt}-free? & Task\hspace*{0pt}-aware? & Eval'd on multi\hspace*{0pt}-fragment? & Deployed in this thesis? \\
\hline
LLMLingua family~\cite{jiang2023llmlingua, jiang2024longllmlingua, pan2024llmlingua2} & yes (inference-only) & question-conditioned & no & LLMLingua-2 \\
Selective Context, RECOMP & yes & yes & no & -- \\
Gist Tokens / AutoCompressor / ICAE~\cite{mu2023gist, chevalier2023autocompressor, ge2024icae} & no (training required) & varies & no & -- \\
Extractive small-LLM (Phi-3 family)~\cite{phi3} & yes (inference-only) & yes & no & Phi-3-Mini ext. \\
Instruction-aware filter & yes & yes & no & yes (this thesis) \\
Truncation (prefix-keep) & yes & no & no & yes (baseline) \\
\hline
\textbf{This thesis (multi-fragment cliff sweep)} & \textbf{yes} & \textbf{yes} & \textbf{yes} & \textbf{four compressors} \\
\hline
\end{tabular}
\end{table}

\section{Multi-Agent LLM Systems and Agentic Memory}
\label{sec:rw_multiagent}

The multi-agent LLM-systems literature has grown rapidly since
2023. The most-cited frameworks are
AutoGen~\cite{wu2023autogen} (an open-source framework for
multi-agent conversation; originally appeared at the ICLR 2024
LLM-Agents Workshop, where it won Best Paper, and was subsequently
published at the Conference on Language Modeling (COLM) 2024),
MetaGPT~\cite{hong2024metagpt} (ICLR 2024 main),
CAMEL~\cite{li2023camel} (NeurIPS 2023),
Reflexion~\cite{shinn2023reflexion} (NeurIPS 2023),
MemGPT~\cite{packer2023memgpt} (arXiv preprint, 2023), and
Collaborative Memory~\cite{rezazadeh2025collaborative} (Rezazadeh et al., 2025 preprint). Each
framework implicitly answers how shared context flows between agents
and what an individual agent needs to see. None of them formally
bounds the information lost when that shared context is compressed.

The agentic-memory line of work
\cite{packer2023memgpt, xu2025amem, chhikara2025mem0,
rasmussen2025zep, saleh2025memindex, saleh2025messagebrokers}
addresses the related question of how to store and retrieve
agent state across long horizons. The closest published academic
analogue to this thesis's memory bus is
MemIndex~\cite{saleh2025memindex}, which the three-layer memory-%
bus architecture of \charef{implementation} explicitly
extends. The MemGPT-style ``LLM as operating system''
framing~\cite{packer2023memgpt} is related but is concerned with
single-agent long-horizon memory rather than multi-agent shared
state. Mem0~\cite{chhikara2025mem0} is a production-oriented
agentic-memory service.

\emph{What is missing.} These frameworks treat compression as an
implicit deployment optimisation; none measures the
\emph{multi-fragment} coordination-quality cost of compression
structurally. This thesis fills that gap (the cliff sweep and the
disclosure measurement); the orthogonal multi-round-agent gap is
left to the AutoGen backend in \texttt{src/m6/orchestrator/} as
future work.

\section{Retrieval-Augmented Generation and Long-Context Benchmarks}
\label{sec:rw_rag}

Retrieval-augmented generation as a paradigm dates to Lewis
\emph{et~al.}~\cite{lewis2020rag}. The landmark works since
include RAPTOR~\cite{sarthi2024raptor} (recursive abstractive
processing for tree-organised retrieval, ICLR 2024),
GraphRAG~\cite{edge2024graphrag} (NeurIPS 2024),
HippoRAG~\cite{gutierrez2024hipporag, gutierrez2025hipporag2}
(NeurIPS 2024 and follow-on), and
Self-RAG~\cite{asai2024selfrag} (ICLR 2024). The pipeline
catalogue this thesis evaluates (P1, P2, P3) places compression
on the retrieval critical path explicitly; the closest published
analogue to the post-retrieval-compression pipeline (P2) is the
LongLLMLingua configuration~\cite{jiang2024longllmlingua}, and
the closest published analogue to the joint
relevance-conditional pipeline (P3) is the adaptive
context-compression framework (ACC-RAG) of Guo
\emph{et~al.}~\cite{guo2025dynamic}, which varies the compression
rate by input complexity.

The long-context-benchmark literature
(\cite{bai2024longbench, hsieh2024ruler, liu2024lostinmiddle,
mialon2023gaia, liu2024agentbench, yen2025helmet}) measures how
effectively a model uses a long context. The
``lost in the middle'' result~\cite{liu2024lostinmiddle} is the
canonical observation that information placed in the middle of a
long context is recovered less reliably than information at the
boundaries. RULER~\cite{hsieh2024ruler} extends the
needle-in-a-haystack methodology to a broader set of long-context
tasks. LongBench~\cite{bai2024longbench} is the standard
multi-task benchmark. None of these benchmarks measures the
multi-fragment-coordination task structure of this thesis;
they measure single-agent long-context comprehension. They are
complementary to the cliff measurement rather than competitive
with it: long-context benchmarks measure the capacity of the
model, while the cliff measurement measures the compression-%
quality effect on a task structure the long-context benchmarks
do not include.

The canonical multi-hop question-answering benchmarks include
HotpotQA~\cite{yang2018hotpotqa} (EMNLP 2018),
2WikiMultiHopQA~\cite{ho2020wiki2hop} (COLING 2020), and
MultiHop-RAG~\cite{tang2024multihoprag} (COLM 2024).
HotpotQA and MultiHop-RAG serve as the two real-data external
arms on which Corollary~2 of \charef{experiments} validates
the information-density-scaling claim; the third data point
in the $n{=}3$ comparison is the synthetic C1 family-a.
The transfer is qualitative: the absolute cliff positions
differ between synthetic and real benchmarks because the
information density $\theta_\text{info}$ differs, but the
cliff structure (existence, sharpness as a function of
$\theta_\text{info}$) transfers.

\emph{What is missing.} RAG benchmarks measure retrieval quality
and long-context benchmarks measure model capacity; neither captures
the multi-fragment compression-quality effect the cliff measurement
isolates. The H3 catalogue fills the
compression-on-the-retrieval-critical-path gap.

\section{Privacy in Agentic Memory and Compressed Retrieval}
\label{sec:rw_privacy}

The privacy-aware-retrieval literature has two strands.
The first protects the retrieval index itself, on the assumption
that direct access to the index would leak protected facts.
Examples include PrivacyRAG~\cite{zhou2025privacyrag} and
SecureRAG~\cite{bassit2025securerag}. The second strand protects
the model weights from training-data extraction, which is a
property of the training pipeline rather than of the deployed
retrieval-and-compression service.

The inference-disclosure metric this thesis contributes (C4) is
distinct from both: it measures leakage \emph{through the
compressor itself}, on the assumption that the retrieval index
is access-controlled but the compressed summaries are exposed to
the planner and to any downstream reader the planner shares them
with. The threat model is intuitive: an adversary who can query
the planner can ask questions about the protected fact, and the
disclosure rate is the fraction of those questions on which the
planner's answer recovers the protected fact given that it has
the compressed summary as evidence. The H4 measurement
operationalises this in a controlled setup; the memory-bus
policy layer of \charef{implementation} operationalises
the resulting disclosure budget as a deployment-time enforcement
mechanism.

The CFedRAG framework~\cite{addison2024cfedrag} is a related
piece of work on coordination of federated retrieval-augmented
generation; it addresses the retrieval-index-level privacy
question rather than the compressor-level disclosure question
that this thesis takes up.

The closest published prior work that connects compression
directly to privacy is
SecurityLingua~\cite{li2025securitylingua} (CoLM 2025), from the
LLMLingua group, which uses prompt compression as a
\emph{defense} against jailbreak attacks: the compressor is
trained to strip adversarial structure from inputs before they
reach the model. SecurityLingua's framing is the dual of
C4's: SecurityLingua applies compression as an input-side
filter against the \emph{adversary's} prompt, whereas C4 measures
compression as an unintended leakage channel for protected facts
in the \emph{defender's} stored content, with a per-question
disclosure rate against a fixed downstream reader. The two
contributions are complementary rather than overlapping (a
deployed memory bus could in principle combine
SecurityLingua-style input filtering with C4-style summary-level
disclosure auditing) but they measure different sides of the
compression-versus-privacy interaction.

\emph{What is missing.} No published work, to the author's
knowledge, measures \emph{summary-level} inference disclosure: the
rate at which a reader recovers a protected fact from a compressed
summary of content it has never seen.
SecurityLingua~\cite{li2025securitylingua} measures the dual
(input-side defense); the privacy-aware-retrieval line protects the
index, not the summary. This thesis fills that gap, and the policy
layer converts the disclosure rate into a deployment budget.

\section{The Gap This Thesis Closes}
\label{sec:rw_gap}

Four sentences, each one mapping to one of the four contributions
of \charef{introduction}.

There is no shared multi-fragment coordination benchmark
probing the cliff in a controlled, reproducible setup. C1 is the
benchmark.

There is no published operational bound that predicts cliff
position from compressor-side and task-side measurements within
a defined regime. The compounding-error model and its
cross-architecture frontier validation are the bound.

There is no honest catalogue of RAG-plus-compression pipeline
placement under matched cost regimes, with the cost model
explicit in EUR per workflow. The C3 catalogue is the
catalogue, and its honest finding is that the assumed
sign-flip does not appear.

There is no published academic measurement of inference
disclosure through the compressor at the summary level. The H4
measurement, integrated with the memory bus's policy layer, is
the measurement.

The thesis sits at the intersection of three threads: the
long-context benchmark suite and the LLMLingua line of work (which
measure orthogonal single-agent properties) and Anthropic's
multi-agent report~\cite{anthropic2025multiagent} (an industry
observation, not a controlled measurement), and contributes the
controlled cross-compressor multi-fragment evaluation none of them
provides: it measures compression's effect on \emph{multi-fragment}
task success, where LLMLingua-2~\cite{pan2024llmlingua2} measures
single-agent QA-F1 retention, LongLLMLingua~\cite{jiang2024longllmlingua}
post-retrieval long-context QA, MemIndex~\cite{saleh2025memindex}
agent-memory storage, and the Anthropic report token-usage-versus-%
variance on one workflow.
